\section{Introduction}\label{sec:intro}

Given a symmetric matrix $G\in\mathbb{R}^{n\times n}$, the nearest correlation
matrix problem asks for
\begin{equation}\label{eq:ncm}
  \min_{X}\ \tfrac12\fro{X-G}^2 \quad\text{subject to}\quad X\succeq 0,\ \diag(X)=e .
\end{equation}
The problem has a unique solution. Since Higham's alternating-projections
method~\cite{higham2002nearest}, the literature has produced dual formulations and
first-order schemes~\cite{malick2004dual,boyd2005least}, a quadratically convergent
semismooth Newton method~\cite{qisun2006quadratically}, its preconditioned and globalized
refinement~\cite{borsdorf2010preconditioned}, accelerated projection
methods~\cite{higham2016anderson}, and, most recently, accelerated gradient
methods with quasi-Newton preconditioning~\cite{huynh2025accelerated}. Most of these
papers include numerical comparisons, and those comparisons are how a reader
decides which method to use.

This paper is about how such comparisons are made. We ask a narrow question:
\emph{when NCM solvers are compared, which methodological choices change the
conclusion, and which only change the numbers?} We identify four candidate
confounds:
\begin{enumerate}
  \item \emph{Native stopping rules.} Each method stops when its own criterion
    is met, so the solvers stop at different accuracies.
  \item \emph{Approximate references.} Forward error is measured against the
    output of another solver instead of against the true solution.
  \item \emph{Rejected-iterate accounting.} Line-search trials that are
    rejected can be credited with reaching a target that the method has not
    actually reached.
  \item \emph{Hidden primitive work.} `Iterations' count different
    operations in different methods, and some per-iteration work is not
    recorded at all.
\end{enumerate}
A careful benchmark must control each of these. Without an exact solution,
however, the second confound cannot be removed, and without trajectories the
first and third cannot be separated.

\paragraph{Contributions.}
In brief, we give an exact-solution NCM family and a matched-accuracy
protocol, and use them to show that conclusions drawn from NCM solver
comparisons do not transfer across benchmark regimes: the rank at which an
ordering established at $n=100$ reverses moves higher at $n=500$, the order
also changes with the degree of invalidity on matrices built from equity
returns, and four matrices from the literature are consistent with this
dependence. We also show that a known
finite-precision line-search defect recurs in a second published method and
occurs on real matrices.
Each contribution builds on established practice (\cref{sec:related}).
\begin{enumerate}
  \item \emph{An exact-solution NCM family} (\cref{sec:family}). We specialize
    the standard construction of test problems from prescribed optimality
    conditions~\cite{lenard1984randomly,calamai1993new} to the NCM problem.
    For semidefinite programs, the same idea is due to Wei and
    Wolkowicz~\cite{wei2010generating}. The parts specific to NCM are the
    projection identity $\Proj(G+\Diag\ys)=\Xs$, independent control of the
    rank, the multiplicity of the smallest eigenvalue of $C(\ys)$ and its
    distance from zero (the strict-complementarity defect equals the
    multiplicity when that distance is zero and vanishes otherwise, so 45 of
    the 270 instances at each dimension are not strictly complementary), a paired-seed
    design, sufficient conditions for $|G_{ij}|\le1$ together with a
    fixed-rank saturation result for $\max_{i\ne j}|\Xs_{ij}|$, and released
    instance sets at $n=100$ and $n=500$.
  \item \emph{A matched-accuracy protocol for NCM} (\cref{sec:protocol}). We
    adapt data profiles~\cite{more2009benchmarking}, fixed-target
    evaluation~\cite{hansen2021coco} and work--precision
    diagrams~\cite{hairer1996solving} to spectral solvers. The unit of cost is
    one \EVD, which we check against wall-clock time. Because the reference is
    exact, two rules become enforceable: rejected trials receive no accuracy
    credit, and accuracy is matched in forward error. Primitive work is
    reported separately.
  \item \emph{An invariant first place above a regime-dependent ordering}
    (\cref{sec:ranking}). Six verified variants of five methods were run on
    270 instances at $n=100$ and 270 at $n=500$, with 135 further instances at
    higher rank, 125 built from two equity panels at $n=550$ and $n=2105$, and
    four matrices from the literature. The cheapest method is the same under
    all 11 conventions examined, and its cost is flat in dimension: 5
    eigendecompositions at $n=100$ and at $n=2105$. Second at $n=100$ is
    Anderson-accelerated alternating projections~\cite{higham2016anderson},
    under 10 of the 11 conventions and tied for second under the eleventh; it
    keeps that place on the exact-solution families but not on perturbed or
    literature matrices.
    Below the first place the ordering changes with the solution rank and
    with how far the input is from a correlation matrix, and, for SBB-Dual
    against AGD-SDAJ-BH, with the accuracy target. The rank dependence holds
    between published methods alone:
    Dykstra-APM and AGD-SDAJ-BH exchange places between $r=20$ and $r=50$ at
    $n=100$, and the exchange has not occurred by $r=150$ at $n=500$. At
    $n=2105$ the order of Anderson-APM and AGD-SDAJ-BH is not monotone in
    rank: Anderson-APM is cheaper at $r=50$ and $r=400$, AGD-SDAJ-BH at $r=100$
    and $r=200$. The same rank pattern appears between SBB-Dual and
    AGD-SDAJ-BH, whose exchange is bracketed between $r=5$ and $r=20$ at
    $n=100$ and between $r=50$ and $r=75$ at $n=500$. Only a family that varies
    the spectral structure at the solution in controlled steps makes such an
    association directly observable.
  \item \emph{Replication of a known line-search defect}
    (\cref{sec:traps}). The failure of objective-value sufficient-decrease
    tests near a minimizer is well known~\cite{hager2005new}. For the NCM dual,
    Borsdorf analysed it~\cite[\S4.7.1]{borsdorf2007newton}, and this led to
    the modified step selection of Borsdorf and
    Higham~\cite{borsdorf2010preconditioned}. The new content is limited to
    two points. First, the defect recurs in a separately published
    first-order method~\cite{huynh2025accelerated}. Second, we quantify it on
    the controlled family: 56/270 stalls for semismooth Newton and 18/270 for
    AGD-SDAJ, both removed by the Borsdorf--Higham correction. We also
    document a second, simpler trap (a redundant Jacobian factorization).
  \item \emph{An NCM quantification of the stopping-rule confound}
    (\cref{sec:stopping}). Dolan, Mor\'e and Munson~\cite{dolan2006optimality}
    warned that native stopping tests confound comparisons. For NCM, the
    widely used rule $\nrm{\nabla\theta}_2\le 10^{-7}n$ leaves the achieved
    forward error spread over a factor of 383 across solvers at $n=100$, and
    it degrades accuracy with dimension unequally across methods.
\end{enumerate}

\paragraph{What we do not claim.}
We do not claim that any method is best in general. The ranking results are
for a synthetic family at $n=100$ and $n=500$, with instances built from
equity returns at $n=550$ and $n=2105$ (\cref{sec:thesis,sec:us2105}) and four
matrices from the literature (\cref{sec:real}) as external checks, all on one
machine. SBB-Dual is a new method from the author's companion manuscript in
preparation~\cite{ibrahim2026sbb}. It is specified in full in
\cref{app:algorithms}, receives the same exposition as the others, and its
convergence theory is not used here. Three findings do not depend on it:
without SBB-Dual, Newton-SIN-BH is still first, Anderson-APM is still
second at $n=100$, and the rank dependence is still present, between
Dykstra-APM and AGD-SDAJ-BH. Two findings do use it: the dependence of the
order on the accuracy target, and the bracketed positions of the rank
crossover at $n=100$ and $n=500$ (\cref{sec:cells,sec:highrank}).

\subsection{Related work}\label{sec:related}

\paragraph{Benchmarking methodology.}
Performance profiles~\cite{dolan2002benchmarking} compare solvers by the ratio
of their cost to the best cost on each problem. Cost is measured at each
solver's own termination, so the solver's stopping test implicitly sets the
accuracy. Dolan, Mor\'e and Munson~\cite{dolan2006optimality} showed that
differing stopping criteria bias such profiles, and proposed rerunning under a
common optimality measure. Data profiles~\cite{more2009benchmarking} count
cost in a problem-independent unit and show how rankings change with the
required accuracy. COCO~\cite{hansen2021coco} uses a ladder of fixed targets
against known optima. Work--precision diagrams, standard in the comparison of
ODE solvers~\cite{hairer1996solving}, plot error against a reference solution
versus work. Gould and Scott~\cite{gould2016note} show that performance-profile
rankings below the leading solver can change when solvers are added or
removed. Beiranvand, Hare and Lucet~\cite{beiranvand2017best} survey best
practice, including known-solution test problems and reproducibility. Our
protocol follows these principles; our contribution is to make their
exact-reference versions possible for the NCM problem. The second/third-place
reversals of \cref{sec:ranking} arise with a \emph{fixed} solver set, from the
accuracy target and the instance regime, so the mechanism differs from the one
studied by Gould and Scott.

\paragraph{Test problems with known solutions.}
Generating test problems from prescribed optimality conditions is standard
for unconstrained optimization~\cite{more1981testing} and for quadratic
programming~\cite{lenard1984randomly,calamai1993new}. Wei and
Wolkowicz~\cite{wei2010generating} generate semidefinite programs from a
chosen optimal primal--dual pair with prescribed complementarity nullity, and
show that nullity correlates with solver difficulty. Our family applies the
same idea to the NCM problem. Their generator does not apply directly: it
builds a linear-objective SDP and chooses the constraint operator as part of
the construction, whereas NCM fixes the constraint operator as $\diag(X)=e$
and leaves only $G$ free. As far as we have found, no NCM study uses
instances with a known exact solution. The test sets
of~\cite{higham2002nearest,qisun2006quadratically,borsdorf2007newton,borsdorf2010preconditioned,higham2016anderson,armijo2025semismooth}
are random or data-derived, as are the Anymatrix \texttt{CORRINV}
matrices~\cite{higham2022anymatrix} used in~\cite{huynh2025accelerated}, and
in all of them the reference solutions are computed approximately.

\paragraph{Finite-precision line searches.}
An objective-value sufficient-decrease test cannot be evaluated accurately near
a minimizer, because the difference of function values suffers cancellation.
Hager and Zhang~\cite{hager2005new} introduced approximate Wolfe conditions
for this reason. For the NCM dual, Borsdorf~\cite[\S4.7.1]{borsdorf2007newton}
derived the stagnation threshold, $\nrm{\nabla\theta}$ of the order of the
square root of the rounding error in $\theta$ (about $10^{-8}$), noted that
the problem is general to Armijo backtracking, and proposed the alternatives
that became the modified step selection of~\cite{borsdorf2010preconditioned}.
\Cref{sec:traps} replicates this known phenomenon; it does not rediscover it.

\subsection{Outline}
\Cref{sec:family} constructs the instance family. \Cref{sec:protocol} states
the protocol, \cref{sec:traps} describes the two implementation traps,
\cref{sec:ranking} reports the ranking study, and \cref{sec:stopping}
examines the dimension-scaled stopping rule. \Cref{sec:repro} covers
reproducibility and \cref{sec:discussion} discusses limitations. The
appendices contain the proofs for the instance family, the verification
records for our reimplementations of AGD-SDAJ and Anderson-APM, the artifact
manifest, and pseudocode for every solver.
